% !TeX root = main.tex
\chapter{Conclusions}
\label{chap:conclusions}

This dissertation developed \texttt{fedctl}, built a Raspberry Pi cluster with \(30\) \texttt{rpi4} and \(20\) \texttt{rpi5} workers, evaluated heterogeneous FL methods under realistic deployment conditions, and proposed \texttt{FedCover} for asynchronous submodel FL. The central argument is that heterogeneous FL methods must be evaluated through both model performance and deployment behaviour, since deployment factors such as device capability, memory constraints, update timing, network conditions, submodel quality, and parameter coverage shape outcomes.

Taken together, the results support three claims. Firstly, compute and straggler-induced heterogeneity create trade-offs between wall-clock time, client participation, submodel quality, and fairness. Secondly, asynchronous submodel FL is feasible without redesigning either underlying method family, but partial, asynchronous replies make parameter coverage an important factor in aggregation. Thirdly, \texttt{FedCover}'s capped inverse-coverage correction improves time-to-target performance relative to uncorrected \texttt{Async-HeteroFL}.

\section{Work completed}

\paragraph{1. Developed \texttt{fedctl}: a tool for deployed FL evaluation.} \hangindent=2em
To fulfil main aim~\hyperref[aim:a]{(a)}, I designed and implemented \texttt{fedctl} as the tool used to run and inspect the experiments in Chapter~\ref{chap:evaluation}. It packages a Flower application, records the effective run and deploy configs, submits the request to a queued service, and launches the corresponding Nomad jobs on the Raspberry Pi cluster. The CLI and browser surfaces expose queue state, logs, artefacts, and results, while the run-config/deploy-config split keeps method choices separate from placement, resources, and network conditions. This made the evaluation reproducible rather than dependent on ad hoc launch scripts.

\paragraph{2. Built a heterogeneous Raspberry Pi cluster for deployed FL evaluation.} \hangindent=2em
To fulfil main aim~\hyperref[aim:b]{(b)}, I built and provisioned a Raspberry Pi cluster with \(30\) \texttt{rpi4} and \(20\) \texttt{rpi5} workers, configured with Ansible and orchestrated through Nomad. \texttt{fedctl} exposes the cluster as typed workers, applies controlled \texttt{netem} impairment, and records device-labelled diagnostics. The testbed turns physical heterogeneity into controlled experimental variables with a traceable evidence path through configuration, logs, W\&B metrics, artefacts, and runtime metadata.

\paragraph{3. Evaluated compute heterogeneity with submodel FL.} \hangindent=2em
To fulfil main aim~\hyperref[aim:c]{(c)}, I implemented and evaluated \texttt{FedAvg}, \texttt{HeteroFL}, \texttt{FedRolex}, and \texttt{FIARSE}, covering aim~\hyperref[aim:c.1]{(c.1)}--\hyperref[aim:c.3]{(c.3)}. Section~\ref{sec:eval_compute_heterogeneity} showed that submodel FL helps when local compute is the bottleneck, but the best method depends on the objective: \texttt{FedAvg} remains the full-model quality control, \texttt{FIARSE} gives the strongest deployable submodel quality, and \texttt{HeteroFL}/\texttt{FedRolex} give the clearest wall-clock savings.

\paragraph{4. Evaluated straggler-induced heterogeneity with asynchronous FL.} \hangindent=2em
To fulfil extension aims~\hyperref[eaim:a]{(a)} and \hyperref[eaim:b]{(b)}, I evaluated straggler-induced heterogeneity using synchronous \texttt{FedAvg} as a reference and asynchronous \texttt{FedAsync}, \texttt{FedBuff}, and \texttt{FedStaleWeight} under mixed completion times and controlled network conditions. Section~\ref{sec:eval_network_heterogeneity} showed that asynchronous updates are useful mainly when client completion times differ: in all-\texttt{rpi5} controls, synchronous \texttt{FedAvg} remains the stronger target-attainment baseline, while mixed \texttt{rpi4}/\texttt{rpi5} runs make \texttt{FedAsync} fastest and \texttt{FedStaleWeight} better at restoring slow-device influence.

\paragraph{5. Combined asynchronous FL with submodel FL and proposed \texttt{FedCover}.} \hangindent=2em
To fulfil extension aims~\hyperref[eaim:c]{(c)} and \hyperref[eaim:d]{(d)}, I evaluated buffered asynchronous aggregation over nested submodel updates and introduced \texttt{FedCover} in Section~\ref{sec:fedcover_impl}. Unlike workload-adaptive prior work \parencite{zhang2023timelyfl}, this dissertation fixes the nested model-rate assignment and isolates aggregation-side coverage correction. Section~\ref{sec:eval_async_submodel_fl} shows that the direct composition, \texttt{Async-HeteroFL}, is feasible, but stale-weighted buffers make parameter coverage dynamic. \texttt{FedCover}'s capped inverse-coverage correction improves time to target while leaving slow-device representation close to the uncorrected baseline.

\section{Future work}
\label{sec:future_work}

\paragraph{Broader workloads, devices, and runtimes.} \hangindent=2em
The evaluated workloads expose compute and runtime trade-offs on CPU-only Raspberry Pi clients, but leave broader tasks and heterogeneous devices open \parencite{caldas2018leaf,lai2022fedscale,li2021ditto}. Future work should test whether the same trade-offs hold on a wider client-device mix, including accelerators such as NVIDIA Jetson Nano devices and Raspberry Pi 5 QPU runtimes.

\paragraph{Representation-aware asynchronous training.} \hangindent=2em
Weighting stale updates can improve slow-device influence, but it cannot make slow clients arrive more often. Device-correlated skew therefore needs methods that combine staleness weighting with fairness-aware objectives, device-aware admission, client selection, or data-aware scheduling \parencite{li2019qffl,nishio2018fedcs,lai2021oort,chai2020tifl}.

\paragraph{Hardware-aware submodel execution.} \hangindent=2em
\texttt{FIARSE} improves submodel quality here, but executes sparse masks through dense PyTorch kernels on edge CPUs. Future work should separate the algorithmic value of sparse importance-selected submodels from the systems cost of executing them without sparse-kernel or hardware support.

\paragraph{Trace-driven straggler-induced heterogeneity evaluation.} \hangindent=2em
The \texttt{netem} make impairment reproducible, but they remain controlled approximations of real deployment conditions. Future work should replay measured edge-network traces through the same \texttt{fedctl} path, connecting controlled profiles with real path variability \parencite{bonawitz2019towards,lai2022fedscale}.

% This dissertation argues that FL systems should be evaluated as deployed distributed systems rather than idealised statistical abstractions.
